\documentclass[12pt]{article}
\usepackage[utf8]{inputenc}
\usepackage[T1]{fontenc}
\usepackage{hyperref}
\title{Wudu Language Resources for K-12 Curriculum Development}
\author{Compiled by Claude Sonnet 5.5 (claude-sonnet-5-5) for Andr\'as Kornai}
\date{}
\begin{document}
\maketitle

\section{Introduction}
Wudu (ISO 639-3 wud, Glottocode wudu1238) is a Gbe language of the Atlantic-Congo family, spoken in Togo near Atakpam\'e (Amou Prefecture, Plateaux Region). Wikipedia gives about 2,500 speakers in 2019. Joshua Project gives roughly 3,000. Glottolog places it in Eastern Gbe (Fongbeic).

\section{Primary Educational Resources}
No primer, textbook or school programme was found, and the language is not known to be taught in schools.

\section{Linguistic and Reference Materials}
\begin{itemize}
\item Wudu language (Wikipedia): \url{https://en.wikipedia.org/wiki/Wudu_language}
\item Glottolog languoid map for Wudu: \url{https://glottolog.org/resource/languoid/id/wudu1238.bigmap.html}
\item ThoT database, tonal status of Wudu: \url{https://thot.huma-num.fr/db/wudu1238/tonal_status}
\end{itemize}
The ThoT entry rests only on a 1983 Togolese government lexical survey. No dictionary or grammar was found.

\section{Religious and Community Resources}
Joshua Project reports that the language is written and that Bible portions exist, with translation work under way. No specific scripture title or URL was seen in the search results, so none is listed.
\begin{itemize}
\item Joshua Project, Wudu language: \url{https://joshuaproject.net/languages/wud}
\item Joshua Project, Wudu people in Togo: \url{https://joshuaproject.net/people_groups/19575}
\end{itemize}

\section{Digital and Audio Resources}
No audio, video or corpus was found.

\section{Curriculum Development Context}
Online material is minimal. The language sits in a dialect continuum with Ewe and Gen, which have more written resources. Any curriculum would first need basic documentation, an orthography and community input.
\end{document}